% LỜI TỰA --- Quyển XIV
\chapter*{Lời Tựa}
\addcontentsline{toc}{chapter}{Lời Tựa}
\markboth{Lời Tựa}{Lời Tựa}

\begin{truyen}{Lời Tựa}{100 Questions Students Ask About Learning}
\chuhoa{G}{iáo viên bắt đầu bộ sách này với một câu hỏi đơn giản: làm thế nào để mỗi buổi học có thể để lại thứ gì đó đáng nhớ hơn là bài thi?}
\textit{(A teacher began this series with a simple question: how can each class leave behind something more memorable than a test?)}

Có những câu hỏi học trò đặt ra thành lời. Và có những câu khác~--- hiện trong ánh mắt, trong cách ngại giơ tay, trong nỗi sợ kiểm tra.
\textit{(Some questions students ask aloud. And others~--- visible in their eyes, in the hesitation to raise a hand, in the anxiety before tests.)}

Quyển mười bốn gom 100 câu hỏi rất gần với đầu óc học sinh. Không nhằm biến việc học thành một bài giảng đạo đức dài, mà để trả lời ngắn, rõ, thật, và đủ gần với những điều các em đang sống trong trường lớp.
\textit{(Volume fourteen gathers 100 questions that stay close to the student mind. Not to turn learning into a long moral lecture, but to answer briefly, clearly, honestly, and close to school life as students actually live it.)}

Bộ sách <<Truyện Tích Cảm Hứng Song Ngữ>> gồm 28 quyển, mỗi quyển một chủ đề riêng. Quyển XIV này mang chủ đề: \textbf{học, động lực, thói quen, áp lực và con đường tương lai}.
\textit{(The series ``Inspiring Bilingual Tales'' contains 28 volumes, each with its own theme. Volume XIV focuses on: \textbf{học, động lực, thói quen, áp lực và con đường tương lai}.})
\end{truyen}

\begin{baihoc}
Học không chỉ là chuyện điểm số~--- học còn là chuyện hiểu mình, hiểu đời, bớt sợ sai, và lớn lên đàng hoàng.
\textit{(Learning is not only about scores~--- it is also about understanding oneself, understanding life, fearing mistakes less, and growing into a grounded person.)}
\end{baihoc}

\begin{ghinhoanh}
\textbf{Short English to remember:}

The right question is already half the answer.

Learning begins the moment you dare to ask why.
\end{ghinhoanh}

\begin{tuhocnhanh}
1. Câu hỏi nào trong 100 câu này em thấy gần với bản thân nhất? \textit{(Which of the 100 questions here feels most personal to you?)}

2. Có câu hỏi nào em muốn thêm vào nếu có một Quyển XV giống quyển này? \textit{(Is there a question you'd add to a future volume like this?)}

3. Điều gì nhất về việc học mà em vẫn chưa tìm được câu trả lời thỏa đáng? \textit{(What about learning do you still not have a satisfying answer for?)}
\end{tuhocnhanh}

\begin{center}
\textit{\color{papergold}``Câu hỏi đúng dẫn tới bài học đúng~--- dù chưa chắc đến câu trả lời dễ.''}
\end{center}

\begin{flushright}
\textbf{Tôi Nguyễn Văn Sang}\\
\textit{GV THPT Nguyễn Hữu Cảnh - TP HCM}
\end{flushright}

\ngancach
